\subsection{Hybrid GDST--Riemann--Siegel Model}

The hybrid Greedy Dirichlet Sub-Sum Transform (GDST)--Riemann--Siegel model provides an exact geometric decomposition of the shifted Riemann zeta function combined with an accelerated functional-equation treatment of the complementary tail. Let \(\alpha_n = \frac{\pi}{n+2}\) for \(n = 0,1,2,\dots\) and let \(\delta_n(\theta) \in \{0,1\}\) denote the additive greedy digits for \(\theta \in [0,\pi]\), i.e., the unique sequence satisfying the greedy selection rule
\[
\sum_{n=0}^{k(\theta)-1} \alpha_n \le \theta < \sum_{n=0}^{k(\theta)} \alpha_n,
\]
where \(k(\theta)\) is the (finite) cutoff index. Define the \emph{selected prefix} (GDST component) by
\begin{equation}
E_{\mathrm{sel}}(\theta,s) := \sum_{n=0}^{k(\theta)-1} \delta_n(\theta) \,(n+2)^{-s}.
\end{equation}
The \emph{rejected tail} is the complementary infinite sum
\begin{equation}
E_{\mathrm{rej}}(\theta,s) := \sum_{n=k(\theta)}^{\infty} (n+2)^{-s}.
\end{equation}
By construction,
\begin{equation}
\zeta(s) - 1 = E_{\mathrm{sel}}(\theta,s) + E_{\mathrm{rej}}(\theta,s), \qquad \operatorname{Re}(s) > 1.
\end{equation}

The \emph{hybrid GDST--Riemann--Siegel model} replaces the raw tail \(E_{\mathrm{rej}}(\theta,s)\) with its Riemann--Siegel accelerated counterpart \(E_{\mathrm{rej}}^{\mathrm{RS}}(\theta,s)\), obtained by applying the Riemann--Siegel approximate functional equation to the tail series beginning at index \(m = k(\theta)+1\). For \(s = \frac12 + it\) with \(t > 0\) sufficiently large, let
\[
N_{\mathrm{tail}}(t;\theta) = \left\lfloor \sqrt{\frac{t}{2\pi}} \right\rfloor - k(\theta),
\]
and define the hybrid tail as
\begin{equation}
\begin{aligned}
E_{\mathrm{rej}}^{\mathrm{RS}}(\theta,s) :={}& \sum_{n=m}^{M} (n+2)^{-s} \\
&+ \chi(s) \sum_{n=1}^{M'} (n+2)^{s-1} \cdot \mathbf{1}_{n \ge m} \\
&+ R_{\mathrm{RS}}(t;\theta),
\end{aligned}
\end{equation}
where:
\begin{itemize}
\item \(M = N_{\mathrm{tail}}(t;\theta)\) is the main-sum cutoff adjusted for the greedy prefix,
\item \(\chi(s) = \pi^{s-1/2} \frac{\Gamma\bigl(\frac{1-s}{2}\bigr)}{\Gamma\bigl(\frac{s}{2}\bigr)}\) is the usual functional-equation factor,
\item \(M'\) is the conjugate-sum cutoff,
\item \(R_{\mathrm{RS}}(t;\theta)\) denotes the explicit Riemann--Siegel remainder series (Edwards 2001; DLMF 25.10), whose leading term satisfies
\[
R_{\mathrm{RS}}(t;\theta) = O\bigl(t^{-1/4}\bigr)
\]
with constants depending on \(\theta\).
\end{itemize}

The full hybrid representation of the zeta function on the critical line is therefore
\begin{equation}
\zeta\!\left(\tfrac12 + it\right) = 1 + E_{\mathrm{sel}}\!\bigl(\theta,\tfrac12 + it\bigr) + E_{\mathrm{rej}}^{\mathrm{RS}}\!\bigl(\theta,\tfrac12 + it\bigr).
\end{equation}

The associated \emph{hybrid rejected phase function} is
\begin{equation}
\psi^{\mathrm{hyb}}(t;\theta) := \arg \bigl( E_{\mathrm{rej}}^{\mathrm{RS}}(\theta,\tfrac12 + it) \bigr),
\end{equation}
with a continuous branch chosen away from the zeros of the tail. The hybrid phase of \(\zeta(\tfrac12 + it) - 1\) decomposes as
\begin{equation}
\arg\bigl(\zeta(\tfrac12 + it) - 1\bigr) = \arg\bigl(E_{\mathrm{sel}}(\theta,\tfrac12 + it)\bigr) + \psi^{\mathrm{hyb}}(t;\theta) + O\bigl(t^{-1/4}\bigr).
\end{equation}

This hybrid model is meromorphic in \(s\) for the selected component (retaining the full pole structure of the original GDST) while the tail admits the standard analytic continuation of the Riemann--Siegel remainder. It therefore furnishes a precise computational and analytic framework for the balance-angle mapping, phase analysis, and self-adjoint spectral realisation pursued in Problems~G1--G6.